\section{An addressed certificate with outward pointers}
\label{sec:construction}

We construct a total function whose accepting certificate belongs to one
column.  The certificate must certify the answers to several Forrelation
instances and point to a zero in every other column.  The latter requirement
will let a genuine accepting column defeat any competitor, even when that
competitor's Forrelation instances violate their promise.

Throughout the construction, fix an integer $k\geq 2$ and let $m$ tend to
infinity through powers of four.  Constants in $O_k(\cdot)$ may depend on
$k$, which remains fixed as $m$ grows.  Write $[m]=\{1,\ldots,m\}$.

\subsection{The partial function and its exact circuit}

We use the $k$-fold Forrelation problem introduced by Aaronson and
Ambainis~\cite{AA18}, with the thresholds analyzed by Bansal and
Sinha~\cite{BS21}. Let $H$ be the normalized Walsh--Hadamard matrix of order $m$, and let
$u=m^{-1/2}(1,\ldots,1)^{\mathsf T}$.  We identify a bit $a$ with the sign
$(-1)^a$.  For a sign vector $z$, write $D_z$ for the diagonal matrix with
diagonal $z$.  Given $k$ sign vectors $z_1,\ldots,z_k\in\{\pm1\}^m$, put
\begin{equation}\label{eq:forrelation-form}
 \Phi(z_1,\ldots,z_k)
 =u^{\mathsf T}D_{z_1}H D_{z_2}H\cdots H D_{z_k}u,
 \qquad \delta=2^{-5k}.
\end{equation}
Define the partial Boolean function $P=P_{k,m}$ on $km$ bits by
\begin{equation}\label{eq:forrelation-promise}
 P(z_1,\ldots,z_k)=
 \begin{cases}
  1,&\Phi(z_1,\ldots,z_k)\geq\delta,\\
  0,&|\Phi(z_1,\ldots,z_k)|\leq\delta/2,
 \end{cases}
\end{equation}
and leave it undefined on all other inputs.  The two displayed cases are
disjoint.  Query complexity for this partial function requires correctness
only on its domain.

\begin{lemma}\label{lem:forrelation}
For each fixed integer $k\geq2$, the functions $P_{k,m}$ have the following
properties.
\begin{enumerate}
 \item For all sufficiently large powers of four $m$,
 \begin{equation}\label{eq:L-def}
  \R(P_{k,m})>L,
  \qquad
  L=\left\lfloor
       \frac{m^{1-1/k}}{(\log_2 m)^2}
     \right\rfloor.
 \end{equation}
 \item For every $0<\eta<1/2$, a quantum algorithm uses
 $O_k(1+\log(1/\eta))$ bit queries and returns $P_{k,m}$ with error at most
 $\eta$ on every input in its domain.  The algorithm can be implemented
 coherently, retaining its workspace.  On every input, promised or not, it
 produces a bit when its output register is measured.
 \item There is a Boolean circuit $C=C_{k,m}$ of fan-in at most two and
 size $S=S_{k,m}=O_k(m(\log m)^2)$ with two output bits $d,b$ such that
 \[
  d=1\quad\Longleftrightarrow\quad P_{k,m}\text{ is defined},
  \qquad d=1\quad\Longrightarrow\quad b=P_{k,m}.
 \]
 The circuit computes these outputs exactly on every input.
\end{enumerate}
\end{lemma}

\begin{proof}
Theorem~1.3 and Corollary~1.4 of Bansal and Sinha~\cite{BS21} apply to
exactly the normalized form and thresholds in
\eqref{eq:forrelation-form}--\eqref{eq:forrelation-promise}.  They give,
for fixed $k$ and error $1/3$,
\[
 \R(P_{k,m})=
 \Omega_k\!\left(
   \left(\frac{m}{\log(km)}\right)^{1-1/k}
 \right).
\]
This exceeds the value of $L$ in \eqref{eq:L-def} for all sufficiently
large $m$.

For the quantum algorithm, prepare $u$, apply the unitary product in
\eqref{eq:forrelation-form}, in its right-to-left order, and test the
projection onto $u$.  Each $D_{z_h}$ uses one phase query, implemented by
the standard bit oracle, while $H$ and the projection require no queries.
The projection succeeds with probability $\Phi^2$.  This probability is
at least $\delta^2$ on a 1-input and at most $\delta^2/4$ on a 0-input.
Independent repetitions and a threshold between these two numbers give
error $\eta$ with $O_k(1+\log(1/\eta))$ repetitions.  The measurement
outcomes and the final threshold decision can instead be computed into
registers and retained, giving the claimed coherent implementation.

For the circuit, let $W=\sqrt m\,H$, whose entries are $\pm1$.  The scalar
\[
 N=z_1^{\mathsf T}W D_{z_2}W\cdots W D_{z_k}(1,\ldots,1)^{\mathsf T}
\]
is an integer, and
\[
 \Phi=\frac{N}{M_{k,m}},
 \qquad M_{k,m}=m(\sqrt m)^{k-1}.
\]
The denominator is an integer because $m$ is a power of four.  Compute
$N$ by $k-1$ unnormalized Walsh--Hadamard transforms, coordinatewise sign
changes, and a final signed sum.  Each transform uses $O(m\log m)$
additions and subtractions by the usual butterfly recursion.  All
intermediate integers have absolute value at most $m^{O(k)}$, so
$O_k(\log m)$ bits suffice.  Binary additions, subtractions, sign changes,
and comparisons have fan-in-two Boolean circuits of size linear in this
bit length.  Consequently the entire arithmetic circuit has Boolean size
$O_k(m(\log m)^2)$.

Finally, the two tests in \eqref{eq:forrelation-promise} are the exact
integer comparisons
\[
 2^{5k}N\geq M_{k,m},
 \qquad 2^{5k+1}|N|\leq M_{k,m}.
\]
Set $d$ to the disjunction of these comparisons and $b$ to the first
comparison.  This adds only $O_k(\log m)$ gates and proves the last
assertion.
\end{proof}

The exact circuit serves a different purpose from the quantum algorithm.
The algorithm will propose an address.  A transcript of the circuit will
certify that this address consists of defined, correct answers, including
on inputs where an arbitrary proposed address has no reason to be correct.

\subsection{The total function}

Set
\begin{equation}\label{eq:r-def}
 r=\lceil20\log_2m\rceil,
 \qquad \ell=\log_2m.
\end{equation}
The input has $m$ columns.  Column $i\in[m]$ contains the following bits:
\begin{enumerate}
 \item A string $X_i\in\bits^m$.
 \item Strings $A_{i,1},\ldots,A_{i,r}\in\bits^{km}$, with no promise
 imposed on them.
 \item A cell $B_{i,s}$ for each address $s\in\bits^r$.  This cell
 contains a pointer $p_{i,s,j}\in[m]$ for each $j\ne i$, encoded by
 $\ell$ bits using a fixed bijection, and an alleged evaluation transcript
 of $C_{k,m}$ on each $A_{i,h}$, $h\in[r]$.
\end{enumerate}
An evaluation transcript consists of one bit for every non-input gate of
$C_{k,m}$, including its output gates.  Circuit inputs are the actual bits
of $A_{i,h}$; they need not be copied into the transcript.  We count these
non-input gates in $S$.

Call column $i$ a \emph{winner} if $X_i=1^m$, every $P(A_{i,h})$ is
defined, and the cell at its true address
\begin{equation}\label{eq:true-address}
 s_i=(P(A_{i,1}),\ldots,P(A_{i,r}))
\end{equation}
contains correct evaluation transcripts and pointers satisfying
\begin{equation}\label{eq:outward-pointers}
 X_j[p_{i,s_i,j}]=0\qquad\text{for every }j\ne i.
\end{equation}
Define $F_{k,m}$ to be 1 if a winner exists, and 0 otherwise.  This is a
total Boolean function: all the preceding conditions are predicates on
the input bits, and no promise is made about the $A$-strings, the pointers,
or the transcripts.

The addressed transcripts have the certificate role used in cheat-sheet
constructions~\cite{ABK16}, while pointers between columns also occur in
earlier query-complexity separations~\cite{GPW15,ABBLSS17}. Here the pointers have
a specific outward orientation: an accepting column supplies a direct
zero witness against every competitor.  In particular, there can be at
most one winner, since a winner's pointer to another column certifies that
the latter is not all ones.  Figure~\ref{fig:outward} depicts this
condition.

\begin{figure}[t]
\centering
\begin{tikzpicture}[>=Latex,font=\small,
 col/.style={draw,rounded corners=2pt,minimum width=2.7cm,minimum height=1.05cm,align=center},
 cell/.style={draw,fill=blue!6,minimum width=3.7cm,minimum height=0.85cm,align=center}]
\node[col] (left) at (-4,1.8) {$X_j$\\$\cdots\ 0\ \cdots$};
\node[col,fill=green!7] (win) at (0,1.8) {$X_{i_\ast}$\\$1\quad1\quad\cdots\quad1$};
\node[col] (right) at (4,1.8) {$X_{j'}$\\$\cdots\ 0\ \cdots$};
\node[align=center,text width=3.7cm,font=\footnotesize] at (0,3.0) {no incoming pointer can find a zero};
\node[cell] (cert) at (0,-0.2) {$B_{i_\ast,s_{i_\ast}}$\\certified address and pointers};
\node[align=center] (addr) at (0,-1.75) {$A_{i_\ast,1},\ldots,A_{i_\ast,r}$\\$s_{i_\ast}=(P(A_{i_\ast,h}))_{h=1}^r$};
\draw[->] (addr) -- (cert);
\draw[->,thick,blue!65!black] (cert.west) to[out=170,in=-90] (left.south);
\draw[->,thick,blue!65!black] (cert.east) to[out=10,in=-90] (right.south);
\draw[dashed] (win.south) -- (cert.north);
\node[align=center,text width=2.8cm,font=\footnotesize] at (-4,-0.85) {zero witness\\disqualifies $j$};
\node[align=center,text width=2.8cm,font=\footnotesize] at (4,-0.85) {zero witness\\disqualifies $j'$};
\end{tikzpicture}
\caption{The relevant cell of a winning column points outward to a zero
in every other column; two such pointers are shown. Only the indicated
zeros are prescribed in the other columns. The cell's address is
computed from the partial-function answers and certified by stored
circuit transcripts. The winner's all-one string prevents every opponent
from disqualifying it, even at an incorrectly decoded address.}
\label{fig:outward}
\end{figure}

\FloatBarrier

We next make certificate verification precise.  For a specified pair
$(i,s)$, say that $i$ wins \emph{at address $s$} if $i$ is a winner and
$s=s_i$.

\Needspace{9\baselineskip}
\begin{lemma}\label{lem:local-tests}
For every input to $F_{k,m}$ and every pair $(i,s)\in[m]\times\bits^r$,
the assertion that $i$ wins at address $s$ is a conjunction of
\[
 K_{k,m}=m+(m-1)+r(S_{k,m}+2)
        =O_k(m(\log m)^3)
\]
Boolean tests.  Each test can be evaluated exactly with $O(\log m)$ bit
queries.  The tests can also be evaluated coherently with the same query
bound, with their indices supplied in superposition.  Their number and
indexing convention are the same for every $(i,s)$.
\end{lemma}

\begin{proof}
Use one test for each bit of $X_i$, requiring it to be 1, and one test
for each $j\ne i$, reading the pointer $p_{i,s,j}$ and requiring the
indicated bit of $X_j$ to be 0.  The latter uses $\ell+1$ queries.

For each of the $r$ transcripts in $B_{i,s}$, use one test per gate,
requiring its claimed value to be the prescribed Boolean function of its
predecessors.  A predecessor value is read from the transcript for a
non-input gate and from $A_{i,h}$ for an input wire.  Fan-in at most two
therefore makes each such test use at most three queries.  Add two output
tests, requiring $d=1$ and $b=s_h$.  By induction in a topological order
of the circuit, all gate tests pass if and only if the transcript is the
true evaluation.  The output tests then say exactly that $P(A_{i,h})$ is
defined and equals $s_h$.  Together with the $X_i$ and pointer tests,
these are precisely the winner conditions at $(i,s)$.

The wiring of the circuit and the locations of all bits are fixed and
require no queries to compute.  A test can load the required bits into
ancillas, compute its result reversibly, and unload the bits by querying
again.  The bit oracle is its own inverse.  A controlled use of that
oracle, when needed, is likewise implemented by loading, controlled use
of the loaded value, and unloading.  Thus coherent evaluation, including
an input-dependent pointer lookup, changes the query count by at most a
constant factor.
\end{proof}

Two consequences will be used repeatedly.  A winning column has one
address at which every test passes; a nonwinning column has a failing
test at every address.  Thus an address proposed on an unpromised input
can always be checked against the total definition.
